% Your notes: each programme, conference and paper below becomes a panel on the
% map, and the same file builds as a PDF. Keep this folder together: main.tex,
% notesmap.sty and img/ (the pictures). Replace the examples with your own;
% everything the notes can say appears once here.
\documentclass[11pt, a4paper]{article}
\usepackage[margin=2.5cm]{geometry}
\usepackage[hidelinks]{hyperref}
\usepackage{notesmap}

\begin{document}
\tableofcontents
\clearpage

% --- a programme: a group's work, pinned to the map by \location -------------
% \begin{programnote}{Title}{key}: the key names it in cross-references (prog:key)
\begin{programnote}{Example Lab -- what the group builds}{Example}
\location{51.50}{-0.12}{London, UK}    % latitude, longitude, the place's name
\thumb{img/example.png}                % the picture on its panel (PNG or JPEG, in img/)

\paragraph{What it is}
One or two sentences on who the group is and what they build. This first
paragraph is the panel's lead line. Their site: \href{https://example.org}{example.org}.

\paragraph{Why it matters here}
Why this programme matters to your work. It builds on the method of the Second
Lab (see \ref{prog:Second}): a cross-reference like that one becomes an arrow
on the map, captioned by its sentence. \aside{A remark for yourself only: it
shows in the PDF, never on the map.}

% a paper inside a programme: listed on the programme's panel
\begin{papernote}{Author et al.\ 2024 -- Short title}{Author2024}
Journal 1:2, \doi{10.0000/example}.    % the paper's link on the map
\thumb{img/paper.png}                  % the picture on the paper's card

\paragraph{Summary}
What the paper did, in plain terms. \gap{Something still to check against the paper.}

\end{papernote}

\end{programnote}

% --- a second programme, so the arrow above has somewhere to go ---------------
% \location[dx,dy]{..}: optionally moves the panel dx,dy pixels away from its pin
\begin{programnote}{Second Lab -- another group}{Second}
\location[300,-200]{-33.87}{151.21}{Sydney, Australia}

\paragraph{What it is}
Another group. A programme with no location waits at the edge of the map
until you give it one.

\end{programnote}

% --- a conference: a programme with dates; the map counts down to its deadline -
\begin{confnote}{Example Conference 2027}{ExConf2027}
\location{48.86}{2.35}{Paris, France}
\dates{2027-06-14}{2027-06-18}         % when it runs (year-month-day)
\deadline[abstracts]{2027-02-01}       % what the map counts down to

\paragraph{What it is}
The meeting, and why you would go. Programme: \url{https://example.org/conf}.

\end{confnote}

% --- a paper on its own, not under any programme -----------------------------
% it sits on the map's shelf of loose papers; cross-reference it as note:Loose2025
\begin{papernote}{Loose et al.\ 2025 -- A paper of no programme yet}{Loose2025}
Preprint, \doi{10.0000/example2}.

\paragraph{Summary}
Read it, then move it into the programme it belongs to.

\end{papernote}

\end{document}
